% 第4章：近似ERM的三步风险比较（高概率形式）。
% 教学数值：参照 w 的总体风险 R(w)=0.10；共同偏差事件内 eta = 0.02；求解容差 alpha = 0.01。
% 最坏情形链：R(h)=0.10 -> Rhat(h) <= 0.12 -> Rhat(hhat) <= 0.13 -> R(hhat) <= 0.15
% 纵坐标：风险 r 映射到 y = 4 + 800 (r - 0.095)（mm）。
\begingroup
\definecolor{FigureBlue}{HTML}{7188A3}\definecolor{FigureBlueFill}{HTML}{E5EAF0}
\definecolor{FigureGreen}{HTML}{829C8B}\definecolor{FigureGreenFill}{HTML}{E8EEE9}
\definecolor{FigureYellow}{HTML}{B4A16F}\definecolor{FigureYellowFill}{HTML}{F1ECDD}
\definecolor{FigureRed}{HTML}{A7787D}\definecolor{FigureRedFill}{HTML}{F0E5E6}
\definecolor{FigureInk}{HTML}{3B4252}\definecolor{FigureMuted}{HTML}{6B7280}
\begin{tikzpicture}[figure picture,x=1mm,y=1mm]
\def\r#1{\fpeval{4+800*(#1-0.095)}}
\node[figure heading] at (22,52) {总体风险};
\node[figure heading] at (62,52) {经验风险};
\draw[draw=FigureGrid,line width=.6pt] (22,2)--(22,49);
\draw[draw=FigureGrid,line width=.6pt] (62,2)--(62,49);
\foreach \v in {0.10,0.12,0.13,0.15}{\draw[draw=FigureGrid,line width=.4pt,dash pattern=on 1pt off 1pt] (6,\r{\v})--(78,\r{\v});
  \node[anchor=east,font=\figurefont\scriptsize,text=FigureMuted] at (6,\r{\v}) {$\v$};}
\node[circle,draw=FigureGreen,fill=FigureGreenFill,line width=.8pt,inner sep=1.3pt] (a) at (22,\r{0.10}) {};
\node[circle,draw=FigureGreen,fill=FigureGreenFill,line width=.8pt,inner sep=1.3pt] (b) at (62,\r{0.12}) {};
\node[circle,draw=FigureBlue,fill=FigureBlueFill,line width=.8pt,inner sep=1.3pt] (c) at (62,\r{0.13}) {};
\node[circle,draw=FigureBlue,fill=FigureBlue,line width=.8pt,inner sep=1.3pt] (d) at (22,\r{0.15}) {};
\node[anchor=east] at (20.5,\r{0.10}) {$R_{\mathcal P}(w)$};
\node[anchor=west] at (63.5,\r{0.12}) {$\hat R_S(w)$};
\node[anchor=west] at (63.5,\r{0.13}) {$\hat R_S(\hat w)$};
\node[anchor=east] at (20.5,\r{0.15}) {$R_{\mathcal P}(\hat w)$};
\draw[figure flow] (a)--node[figure edge label,pos=.5,below right=-.5mm] {$+\eta$} (b);
\draw[figure flow] (b)--node[figure edge label,left] {$+\alpha$} (c);
\draw[figure flow] (c)--node[figure edge label,pos=.5,above right=-.5mm] {$+\eta$} (d);
\node[figure note,anchor=west] at (82,\r{0.125}) {共同偏差事件\\内的最坏情形};
\end{tikzpicture}
\endgroup
